%%%%%%%%%%%%
%
% $Autor: Wings $
% $Datum: 2019-03-05 08:03:15Z $
% $Pfad: Troubleshooting.tex $
% $Version: 4250 $
% !TeX spellcheck = en_GB/de_DE
% !TeX encoding = utf8
% !TeX root = manual 
% !TeX TXS-program:bibliography = txs:///biber
%
%%%%%%%%%%%%

\chapter{Troubleshooting}
Dieses Kapitel  dient dazu, Anwendern bei der schnellen Identifikation und Behebung von Problemen zu helfen. Es bietet eine strukturierte Übersicht über mögliche Fehlerursachen sowie konkrete Lösungsschritte, die leicht nachvollzogen werden können. Ziel ist es, Störungen im Betrieb zu minimieren und eine eigenständige Problemlösung zu ermöglichen. Die beschriebenen Hinweise und Maßnahmen basieren auf typischen Anwendungsszenarien und unterstützen dabei, die Funktionsfähigkeit des Systems effizient wiederherzustellen.
Zuerst werden die Mechanischen Störungen behandelt, anschließend die elektrischen. 

\subsection{Mechanische Fehler}
\begin{itemize}
	\item \textbf{Das Tischförderband lässt sich nicht aufstellen} \\ 
	\textbf{Mögliche Ursache:} Die Klappbaren Füße sind nur blockiert, Es fehlen die Stützbeine \\
	\textbf{Behebung:} Stützbeine auf Fremdkörper prüfen und diese entfernen, die Vorspannung der Beine durch leichtes Lösen der Schraube verringern. 
	\item \textbf{Das Transportband lässt sich nicht bewegen}\\
	\textbf{Mögliche Ursachen:} Fremdkörper blockieren das Transportband \\
	\textbf{Behebung:} Das Förderband auf eingeklemmte Fremdkörper überprüfen. Diese Fremdkörper entfernen. 
	\item \textbf{Das Tischförderband steht schief} \\ 
	\textbf{Mögliche Ursache:} Die Stützbeine des Förderbandes sind nicht vollständig ausgeklappt und arretiert. \\
	\textbf{Behebung:} Die Stützbeine erneut einklappen und ausklappen, die Arretierung der Beine überprüfen.  
	\item \textbf{Das Tischförderband wackelt}\\ 
	\textbf{Mögliche Ursachen:} Das Tischförderband steht auf nicht ebenem Untergrund, das Tischförderband hat ein beschädigtes Stützbein. \\ 
	\textbf{Behebung:} Die Stützbeine des Förderbandes auf Beschädigung überprüfen und ggf. austauschen. Bei keiner Beschädigung das Förderband an einem anderen Ort positionieren und orientieren. 
	\item \textbf{Der Riemen des Förderbandes bewegt sich nicht}\\ 
	\textbf{Mögliche Ursachen:} Das Transportband des Tischförderbandes ist nicht gespannt. \\ 
	\textbf{Behebung:} Das Transportband mittels Spannvorrichtung spannen. 
	
		
\end{itemize}
\subsection{Elektrische Fehler}
Hier werden alle Fehler die direkten Zusammenhang zur Strom- und Spannungsversorgung haben. Zusätzlich wird hier Hilfestellung zu allen Fehlermeldungen auf dem Bediener
\begin{itemize}
	
	\item \textbf{Das Tischförderband startet nicht} \\
	\textbf{Mögliche Ursachen:} Keine Spannungsversorgung, fehlerhafte Verkabelung, defekter Motor, falsche Pin-Konfiguration \\
	\textbf{Behebung:} Spannungsversorgung prüfen, korrekte Schaltposition des Hauptschalters prüfen.
	\item \textbf{das Transportband läuft unregelmäßig schnell} \\
	\textbf{Mögliche Ursachen:} Instabile Spannungsversorgung, PWM-Signal fehlerhaft, mechanische Blockade \\
	\textbf{Behebung:}  Spannungsstabilität sicherstellen, Hersteller kontaktieren. 
	\item \textbf{Fehler: \glqq Sensor liefert keine Daten"} \\
	\textbf{Mögliche Ursachen:} Falsche Initialisierung, defekter Sensor, Kommunikationsfehler (I2C/SPI) \\
	\textbf{Behebung:}  Das elektrische Tischförderband an und abschalten. Das Tischförderband mittels Tastenkombination zurücksetzen. Gegebenenfalls den Sensor austauschen lassen. 
	\item \textbf{ Anzeige reagiert nicht} \\
	\textbf{Mögliche Ursachen:} Softwareabsturz, fehlerhafte Stromversorgung \\
	\textbf{Behebung:}  Reset durchführen, Spannungsversorgung prüfen
	\item \textbf{Fehlermeldung: \glqq Motor überhitzt"} \\
	\textbf{Mögliche Ursachen:} Überlast, Dauerbetrieb, unzureichende Kühlung \\
	\textbf{Behebung:}  Dauerlaufzeit reduzieren, Pausen einbauen, Umgebungstemperatur senken. 
	\item \textbf{Förderband stoppt unerwartet} \\
	Mögliche Ursachen: Sicherheitsabschaltung, Spannungsabfall, Softwarefehler \\
	\textbf{Behebung:}  Spannungsversorgung prüfen, elektrisches Tischförderband zurücksetzen, Hersteller kontaktieren.
	\item \textbf{Fehler: \glqq Kommunikationsprobleme zwischen Modulen"} \\
	\textbf{Mögliche Ursachen:} Bus-Konflikte, falsche Adressen, Leitungsstörungen \\
	\textbf{Behebung:}  Den Hersteller kontaktieren.  
\end{itemize}
Alle auf der Anzeige dargestellten Fehler und Fehlermeldungen könnnen gemäß des Benutzerhandbuches behoben und anschließend quitiert werden. Hierfür muss der Fehler behoben, und anschließend der dafür vorgesehene Taster am Bedienelement betätigt werden. Sollte der Fehler unmittelbar nach der Quittierung erneut auftreten, so ist die Spannungsversorgung zu trennen und der Hersteller zu kontaktieren.
\section*{Normative Referenzen}
Folgende sind die zur Beschreibung der Fehlerbehebung verwendeten Normen.
\begin{itemize}
	\item DIN EN 60204-1: Sicherheit von Maschinen – Elektrische Ausrüstung
	\item DIN EN ISO 13850: Sicherheit von Maschinen – Not-Halt-Funktion
	\item DIN EN 61000: Elektromagnetische Verträglichkeit (EMV)
	\item DIN EN 60034: Elektrische Maschinen
	\item VDE 0100: Errichten von Niederspannungsanlagen
	\item ISO/IEC 61131-2: Speicherprogrammierbare Steuerungen
\end{itemize}







